% lattice.tex — Paper 6, Section II: The Self-Modeling Lattice
% ASSERT_CONVENTION: natural_units=natural, metric_signature=mostly_plus, entropy_base=nats, modular_hamiltonian=K_minus_ln_rho, coupling_convention=H_sum_hxy, state_normalization=Tr_rho_1
\section{The Self-Modeling Lattice}
\label{sec:lattice}

\subsection{Self-modeling forces quantum mechanics}
\label{sec:paper5-summary}

Paper~5~\cite{Paper5} established that faithful self-modeling---the
condition that a physical system $B$ contains a subsystem $M$ that
tracks the state of $B$ through their shared boundary---forces the
local state space to be $M_n(\mathbb{C})^{sa}$, the self-adjoint part
of the $n \times n$ complex matrix algebra, equipped with the
L\"uders sequential product
\begin{equation}
  a \mathbin{\&} b = a^{1/2}\, b\, a^{1/2}.
  \label{eq:luders}
\end{equation}
The real and quaternionic alternatives are excluded by local
tomography: only complex quantum mechanics has
$\dim(V_{BM}) = \dim(V_B) \cdot \dim(V_M)$.  The sequential product
is the unique product compatible with axioms S1--S7 of
Ref.~\cite{vandeWetering2019}, and it inherits $U(n)$ covariance from
the self-modeling structure.  We do not re-derive these results here;
they serve as the starting point for the present paper.

\subsection{Lattice definition}
\label{sec:lattice-def}

We place self-modeling systems on a lattice described by an undirected
graph $G = (V, E)$, where $V$ is a countable set of sites and
$E \subseteq \binom{V}{2}$ is the set of edges encoding nearest-neighbor
adjacency.  At each site $x \in V$ the local algebra is
\begin{equation}
  A_x = M_n(\mathbb{C}),
  \label{eq:local-algebra}
\end{equation}
and for a finite region $\Lambda \subset V$ the regional algebra is
\begin{equation}
  A_\Lambda = \bigotimes_{x \in \Lambda} A_x \cong M_{n^{|\Lambda|}}(\mathbb{C}).
  \label{eq:regional-algebra}
\end{equation}
The quasi-local algebra
$\mathfrak{A} = \overline{\bigcup_{|\Lambda|<\infty} A_\Lambda}^{\|\cdot\|}$
is a uniformly hyperfinite (UHF) C*-algebra of type~$n^\infty$.

\emph{Background dependence.}---The graph topology $G = (V, E)$ is an
irreducible \emph{input} to this construction.  Self-modeling
determines the metric (encoded in the coupling strengths and the
resulting entanglement structure) on a given topology, but not the
topology itself.  This is analogous to general relativity, where
Einstein's equations determine the metric on a given differentiable
manifold without determining its topology.  The spatial dimension~$d$
is set by the graph structure: a $d$-dimensional regular lattice
yields $d$ emergent spatial dimensions.

\subsection{Hamiltonian construction}
\label{sec:hamiltonian}

The interaction between adjacent sites is encoded in a two-site
Hamiltonian $h_{xy} \in A_{\{x,y\}}^{sa}$ for each edge
$\{x, y\} \in E$.  Self-modeling locality---the constraint that model
and body interact only through their shared boundary---implies that
$h_{xy} = 0$ for non-adjacent pairs ($\{x,y\} \notin E$), mapping
self-modeling locality directly to nearest-neighbor Hamiltonian
locality~\cite{Paper5}.

The L\"uders sequential product~\eqref{eq:luders} is covariant under
the adjoint action of $U(n)$:
\begin{equation}
  U(a \mathbin{\&} c)U^\dagger = (UaU^\dagger) \mathbin{\&} (UcU^\dagger)
  \qquad \forall\, U \in U(n).
  \label{eq:covariance}
\end{equation}
The composite sequential product
$(a \otimes b) \mathbin{\&} (c \otimes d) = (a \mathbin{\&} c) \otimes
(b \mathbin{\&} d)$ is covariant under $U(n) \times U(n)$ (independent
rotations at each site).  The coupling, however, physically connects
the two sites through a shared boundary; basis independence requires
that the same rotation be applied to both sites.  This constrains
$h_{xy}$ to be invariant under the \emph{diagonal} $U(n)$:
\begin{equation}
  (U \otimes U)\, h_{xy}\, (U \otimes U)^\dagger = h_{xy}
  \qquad \forall\, U \in U(n).
  \label{eq:diagonal-invariance}
\end{equation}
Full $U(n) \times U(n)$ invariance would kill all coupling, leaving
$h_{xy} \propto \1 \otimes \1$.

By \textbf{Schur--Weyl duality} for the symmetric group $S_2$ acting
on $(\mathbb{C}^n)^{\otimes 2}$, this space decomposes under the
diagonal $U(n)$ action into two irreducible representations:
$\mathrm{Sym}^2(\mathbb{C}^n) \oplus \wedge^2(\mathbb{C}^n)$.  By
Schur's lemma the commutant of $\{U \otimes U\}$ is
$\mathrm{span}\{\1, F\}$, where $F$ is the SWAP operator defined by
\begin{equation}
  \swapop \ket{i}\ket{j} = \ket{j}\ket{i}.
  \label{eq:swap-def}
\end{equation}
Therefore the most general diagonal-$U(n)$-invariant self-adjoint
interaction is
\begin{equation}
  h_{xy} = \alpha\, \1 \otimes \1 + J\, \swapop,
  \label{eq:h-general}
\end{equation}
with $\alpha, J \in \mathbb{R}$.  The identity term shifts the energy
zero; setting $\alpha = 0$, the total lattice Hamiltonian is
\begin{equation}
  \boxed{H = \sum_{\langle x,y \rangle \in E} J\, \swapop.}
  \label{eq:hamiltonian}
\end{equation}

The interaction is \emph{forced}: the form $h_{xy} = J\,\swapop$ is
the unique diagonal-$U(n)$-invariant coupling (up to the overall
energy scale~$J$).  Any anisotropy would break the $U(n)$ covariance
inherited from the L\"uders product.  For $n = 2$, the SWAP operator
takes the familiar form
$\swapop = \frac{1}{2}(\1 \otimes \1 + \vec{\sigma} \cdot \vec{\sigma})$,
and the Hamiltonian~\eqref{eq:hamiltonian} reduces to the isotropic
Heisenberg model.

The sign of~$J$ is \emph{not} determined by the self-modeling
constraints: $J > 0$ gives the antiferromagnetic (AFM) Heisenberg
model, $J < 0$ gives the ferromagnetic (FM) model.  Both signs are
compatible with the algebraic structure.  However, the FM ground state
($J < 0$) is a product state with $\ent(A) = 0$ for all subsystem
sizes, producing no entanglement structure and hence no emergent
geometry.  Only the AFM case ($J > 0$) supports the nontrivial
entanglement that drives the Jacobson argument.  We therefore work
with $J > 0$ throughout; this is an input to the derivation, not
forced by self-modeling.

\subsection{Lieb--Robinson bounds and emergent causal structure}
\label{sec:lr-bounds}

The nearest-neighbor Hamiltonian~\eqref{eq:hamiltonian} with finite
coupling $\|h_{xy}\| = |J|$ satisfies the hypotheses of the
Lieb--Robinson theorem~\cite{LiebRobinson1972, NachtergaeleSims2006}.
For any local observables $A_x$ and $B_y$ supported at sites $x$
and~$y$ separated by graph distance $d(x,y)$, the Heisenberg-picture
commutator is bounded by
\begin{equation}
  \norm{[\tau_t(A_x),\, B_y]}
  \leq C_{\mathrm{LR}}\, \norm{A}\, \norm{B}\,
  e^{-\mu(d(x,y) - \vlr |t|)},
  \label{eq:lr-bound}
\end{equation}
where $\tau_t(A) = e^{iHt} A\, e^{-iHt}$ is the Heisenberg evolution,
$\mu > 0$ is a spatial decay constant, $C_{\mathrm{LR}}$ is a
dimension-dependent prefactor, and $\vlr$ is the \textbf{Lieb--Robinson
velocity}---the maximum speed of information propagation on the
lattice.

For the self-modeling Hamiltonian on a lattice with coordination
number~$z$, the Nachtergaele--Sims
framework~\cite{NachtergaeleSims2006} with exponential decay function
$F_a(r) = e^{-ar}$ at $a = 1$ gives
\begin{equation}
  \vlr = \frac{4ze|J|}{e - 1}.
  \label{eq:vlr}
\end{equation}
For the one-dimensional chain ($z = 2$), this is $\vlr = 8e|J|/(e-1)
\approx 12.66\,|J|$.  Crucially, $\vlr$ is \emph{independent of the
local dimension~$n$}: the SWAP operator $F$ satisfies $\|F\| = 1$ for
all $n \geq 2$, so the Lieb--Robinson velocity depends only on the
coupling strength~$|J|$ and the lattice coordination number~$z$.

The bound~\eqref{eq:lr-bound} defines an effective \emph{light cone}
on the lattice: correlations outside the cone ($d(x,y) > \vlr |t|$)
are exponentially suppressed.  In the continuum limit
(Sec.~\ref{sec:einstein}), $\vlr$ is identified with the emergent
speed of light~$c$.
